\clearpage\begingroup\SourceLeaf{615}
\ArithmeticHead{60}{L}{PHILIPP MAENNCHEN, GAUSS ALS ZAHLENRECHNER.}
\renewcommand{\thefootnote}{\arabic{footnote})}\setcounter{footnote}{0}
\noindent\hspace*{\parindent}Man erkennt, dass
\[
\begin{gathered}
d_{11}=d_0-1\\
d_{12}=d_1-1\\
\mathord{.}\quad\mathord{.}\quad\mathord{.}\quad\mathord{.}\quad\mathord{.}\quad\mathord{.}\quad\mathord{.}\\
d_{18}=d_7-1
\end{gathered}
\]
\noindent\hspace*{3em}allgemein\quad $d_{\nu+11}=d_\nu-1$ ist.

Daraus ergibt sich, dass $d=18$, der einzige Wert, der zweimal auftreten
kann (einmal normaler Weise, und das andere Mal durch den um 1 heruntergedrückten Wert 19), nur für $a>10$ auftreten darf, damit das um 11
Nummern zurückliegende $d$ den Wert 19 hat und somit auf 18 herabgedrückt
werden müsste\footnote{Die erste Ausnahme heisst nämlich nach den gregorianischen Bestimmungen: Der Frühlingsvollmond darf höchstens auf den 18. April fallen. Führt daher die Rechnung auf den 19. April, so muss man
um einen Tag zurückdatieren.}.

Wir kommen somit zu der arithmetischen Festsetzung: Ergibt die Rechnung den 25. April als Osterdatum, und ist dabei $a>10$ und $d=28$ (also
$e=6$), so wird der 18. April genommen.

\textsc{Gauss} gab dieser Festsetzung eine etwas andere Prägung. Wir finden
sie bereits in seiner ersten Veröffentlichung, und dort (Werke VI, S. 79)
lautet sie folgendermassen:

»Gibt die Rechnung $d=28$, $e=6$, und kommt noch die Bedingung
hinzu, dass $11M+11$ mit 30 dividiert einen Rest gibt, der kleiner als 19
ist, so fällt Ostern nicht, wie aus der Rechnung folgt, auf den 25., sondern
auf den 18. April.«

Warum wählt \textsc{Gauss} diese eigentümliche Fassung, während doch die
oben entwickelte erstens viel näher liegt und zweitens auch praktischer in
der Anwendung ist? Es scheint, dass \textsc{Gauss} damit beabsichtigte, die Jahrhunderte auszusieben, für die der zweite Ausnahmefall überhaupt nicht in
die Erscheinung tritt. An der vorhin zitierten Stelle heisst es nämlich weiter:
»Man überzeugt sich leicht, dass dieser Fall nur in denjenigen Jahrhunderten
eintreten könne, da $M$ einen von folgenden acht Werten hat:
\[
2,\ 5,\ 10,\ 13,\ 16,\ 21,\ 24,\ 29.\text{«}
\]

Wir wollen zunächst noch zeigen, dass $d=28$ und $a>10$ identisch ist
mit $d=28$ und $11M+11<19\; (\mathrm{mod}\ 30)$.
\clearpage\SourceLeaf{616}\ArithmeticHead{61}{R}{X. CHRONOLOGISCHE ARBEITEN.}
\setcounter{footnote}{0}
\noindent\hspace*{\parindent}Nach Definition ist
\[
d\equiv19a+M\; (\mathrm{mod}\ 30),
\]
\noindent und da $d=28$ sein soll, so ist
\[
28\equiv19a+M\; (\mathrm{mod}\ 30)
\]
\noindent oder
\[
M\equiv28-19a\equiv28+11a\; (\mathrm{mod}\ 30).
\]
\noindent Daraus folgt:
\[
11M+11\equiv19+a\; (\mathrm{mod}\ 30).
\]
\noindent Durchläuft nun $a$ die Werte von 0 bis 10, so durchläuft der 30er-Rest von
$11M+11$ die Werte von 19 bis 29. Für alle übrigen Werte von $a$ ist dieser
Rest also kleiner als 19, und damit ist der Zusammenhang bewiesen.

Die zu
\[
a=11,\ 12,\ 13,\ 14,\ 15,\ 16,\ 17,\ 18
\]
\noindent gehörigen Werte von $M$ sind
\[
M=29,\ 10,\ 21,\ 2,\ 13,\ 24,\ 5,\ 16,
\]
\noindent und diese hat \textsc{Gauss} in der obigen Darstellung nach der Grösse geordnet.
Mit Hilfe der Formel
\[
M\equiv15+k-p-q\; (\mathrm{mod}\ 30)
\]
\noindent kann man nun leicht das oben angedeutete Aussieben vornehmen. Es ergibt
sich, dass der zweite Ausnahmefall für $k=19,\ 20,\ 21$ eintreten kann, dann
erst wieder für $31,\ 32,\ 33$; dagegen nicht für $k=16,\ 17,\ 18$ und auch nicht
für $k=22,\ .\ .\ .,\ 30$.

Während \textsc{Gauss} bei seiner \textls[120]{ersten} Veröffentlichung die beiden Ausnahmefälle durchaus richtig darstellt, hat er merkwürdigerweise in einer späteren
Darstellung\footnote{Astron. Jahrb. f. d. J. 1814, Berlin 1811, S. 273, Werke XI, 1, S. 199.} den zweiten Fall zwar auf eine einfachere, aber nicht auf die
richtige Form gebracht. Er sagt:

Wenn im gregorianischen Kalender die Rechnung Ostern am 25. April
gibt, setzt man allemal den 18\footnote{Denselben Fehler macht auch \textsc{Delambre}, vergl. die \textit{Bemerkung} von A. \textsc{Loewy} Werke XI, 1, S. 200.}.

Der Fall, dass Ostern auf den 25. April fällt, ereignet sich im XIX. Jahrhundert nur einmal, im Jahre 1886, und im XX. auch nur einmal, 1943.
\clearpage\SourceLeaf{617}\ArithmeticHead{62}{L}{PHILIPP MAENNCHEN, GAUSS ALS ZAHLENRECHNER.}
\setcounter{footnote}{0}
\noindent Es ist möglich, dass \textsc{Gauss} eine flüchtige Rechnung anstellte und in diesen
beiden Jahrhunderten kein Jahr fand, in dem Ostern auf den 25. April fiel,
so dass er die einfachere Regel aufzustellen wagte, zumal er in dem erwähnten
Artikel »eine leichte Methode« geben wollte, »das Osterdatum zu bestimmen«.

Nachdem \textsc{Gauss} dieses chronologische Problem »elegant« gelöst hatte,
reizte ihn ein ganz ähnliches, nämlich die rein zahlentheoretische Bestimmung
des Datums für das jüdische Passahfest. Die Lösung gelang ihm etwa ein
Jahr später. Wir finden darüber die Tagebuchnotiz (No. 117) vom 1. April
1801\footnote{Werke X, 1, S. 560.}.

»Iisdem diebus Pascha Iudaeorum per methodum novam determinare docuimus.«

Diese \textls[120]{Passahformel} veröffentlichte er ebenfalls in der Monatlichen Correspondenz, und zwar in der Mainummer 1802\footnote{Werke VI, S. 80.}. Auch in bezug auf diese
bemerkt er ausdrücklich\footnote{Werke VI, S. 86.}, dass damit

»eine rein arithmetische Regel für die Berechnung des jüdischen Osterfestes gegeben ist, deren Anwendung von aller weiteren Bekanntschaft mit
der Einrichtung des jüdischen Kalenders unabhängig ist.«

Über die Berechnung des Passahdatums findet sich im Nachlass nichts,
wohl aber enthält dieser\footnote{Werke XI, 1, S. 215.} kurze Angaben über die chronologisch wichtigere
Berechnung des »Neumonds Tišri« für jedes jüdische Jahr $A$, und dieses
Datum ist zugleich die Grundlage für das vorausgehende Passahdatum.

Die erwähnten Angaben lassen erkennen, dass \textsc{Gauss} vom »Molad Tohu«,
d. h. vom Neumond Tišri des Jahres 1 der jüdischen Zeitrechnung ausging
und die Anzahl der seit diesem Zeitpunkt bis zum Neumond Tišri des Jahres
$A$ verflossenen Tage bestimmte. Dabei ist die einzige zahlentheoretisch interessante Frage die nach der Anzahl der in diesen $A$ Jahren enthaltenen Schaltmonate. Bei der Bestimmung dieser Anzahl geht \textsc{Gauss} von der Beobachtung
eines Verteilungsgesetzes der 7 Schaltmonate auf die 19 Jahre eines Metonischen
Zyklus aus. Er findet, dass in den 19 aufeinanderfolgenden Gruppen von je
8 Jahren (1 bis 8; 9 bis 16; 17 bis 5; 6 bis 13; 13 bis 1; 2 bis 9 u. s. f.)
\clearpage\SourceLeaf{618}\ArithmeticHead{63}{R}{XI. SCHLUSSBETRACHTUNG.}
\setcounter{footnote}{0}
\noindent nur eine Gruppe (9—16) 2 Schaltjahre enthält, alle übrigen aber 3 Schaltjahre. Wie er daraus vermutlich die gesuchte Anzahl hergeleitet hat, darüber
belehren uns die Bemerkungen von \textsc{Loewy} zu den nachgelassenen Bruchstücken
der Berechnung\footnote{Werke XI, 1, S. 216 ff.}.
\clearpage\endgroup

